% !TEX root = ../main.tex
\section{Ejercicio propuesto 2: aproximación de Knapsack mediante escalamiento}

\subsection{Programación Dinámica basada en valores y esquema FPTAS}
A diferencia de la formulación clásica basada en capacidad $W$, la formulación orientada a valores define $DP(i, v)$ como el peso mínimo necesario para acumular exactamente el valor $v$ utilizando un subconjunto admisible de los primeros $i$ objetos:
\begin{equation}
    DP(i, v) = \min \left\{ \sum_{k \in S} w_k \;\middle|\; S \subseteq \{1, \dots, i\}, \ \sum_{k \in S} v_k = v \right\}
\end{equation}
Con transición:
\begin{equation}
    DP(i, v) = \min\left( DP(i - 1, v), \quad DP(i - 1, v - v_i) + w_i \right)
\end{equation}

El algoritmo FPTAS procede en 6 pasos estructurados:
\begin{enumerate}
    \item \textbf{Filtrado:} Descartar objetos con peso $w_i > W$. Si no queda ningún objeto, retornar solución vacía (cuidado numérico).
    \item \textbf{Factor de escala:} Calcular $V_{\max} = \max_i v_i$ y $K = \frac{\varepsilon V_{\max}}{n}$.
    \item \textbf{Discretización:} Escalar los valores mediante $v'_i = \lfloor v_i / K \rfloor$.
    \item \textbf{DP sobre valores reducidos:} Resolver la tabla dinámica sobre $v'_i$ y pesos $w_i$.
    \item \textbf{Reconstrucción:} Recuperar los índices de los objetos óptimos de la tabla escalada.
    \item \textbf{Evaluación real:} Calcular la suma con valores originales $\text{ALG}_\varepsilon = \sum_{i \in S} v_i$ y verificar peso $\le W$.
\end{enumerate}

\subsection{Código fuente de las implementaciones}

\begin{lstlisting}[language=Python, caption={Programación Dinámica exacta por valores}, label={lst:knap_dp}]
def knapsack_dp_valores(pesos, valores, W):
    n = len(pesos)
    if n == 0 or W <= 0: return 0, 0, [], 0
    V_sum = sum(valores)
    if V_sum == 0: return 0, 0, [], 0
    
    INF = float('inf')
    dp = [[INF] * (V_sum + 1) for _ in range(n + 1)]
    dp[0][0] = 0
    estados = (n + 1) * (V_sum + 1)
    
    for i in range(1, n + 1):
        w, v = pesos[i - 1], valores[i - 1]
        for val in range(V_sum + 1):
            dp[i][val] = dp[i - 1][val]
            if val >= v and dp[i - 1][val - v] + w < dp[i][val]:
                dp[i][val] = dp[i - 1][val - v] + w
                
    mejor_v = max(val for val in range(V_sum + 1) if dp[n][val] <= W)
    curr_v = mejor_v
    seleccion = []
    for i in range(n, 0, -1):
        w, v = pesos[i - 1], valores[i - 1]
        if curr_v >= v and dp[i][curr_v] == dp[i - 1][curr_v - v] + w:
            seleccion.append(i - 1)
            curr_v -= v
    seleccion.reverse()
    return mejor_v, sum(pesos[i] for i in seleccion), seleccion, estados
\end{lstlisting}

\begin{lstlisting}[language=Python, caption={Esquema FPTAS por escalamiento}, label={lst:knap_fptas}]
def knapsack_fptas(pesos_orig, valores_orig, W, epsilon):
    validos = [(p, v, idx) for idx, (p, v) in enumerate(zip(pesos_orig, valores_orig)) if p <= W]
    if not validos or W <= 0: return 0, 0, [], 0, 0.0
    pesos = [item[0] for item in validos]
    valores = [item[1] for item in validos]
    indices_orig = [item[2] for item in validos]
    n = len(pesos)
    
    V_max = max(valores)
    K = (epsilon * V_max) / n
    v_prime = [int(math.floor(v / K)) for v in valores] if K > 0 else list(valores)
    
    mejor_v_prime, _, sel_local, estados = knapsack_dp_valores(pesos, v_prime, W)
    seleccion_orig = [indices_orig[i] for i in sel_local]
    valor_real = sum(valores_orig[i] for i in seleccion_orig)
    peso_real = sum(pesos_orig[i] for i in seleccion_orig)
    assert peso_real <= W
    return valor_real, peso_real, seleccion_orig, estados, K
\end{lstlisting}

\subsection{Resultados experimentales}
Se evaluaron \textbf{20 instancias pequeñas} bajo cuatro valores de tolerancia $\varepsilon \in \{0.50, 0.25, 0.10, 0.05\}$, acumulando 80 evaluaciones. La Tabla \ref{tab:resumen_knap} presenta una muestra representativa:

\begin{table}[H]
    \caption{Evaluación experimental de Knapsack FPTAS en instancias representativas}
    \label{tab:resumen_knap}
    \centering
    \scriptsize
    \begin{tabular*}{\textwidth}{@{\extracolsep{\fill}}lcccccccccc}
        \toprule
        \textbf{Instancia} & $\mathbf{n}$ & $\mathbf{W}$ & $\mathbf{\varepsilon}$ & \textbf{OPT} & $\mathbf{ALG_\varepsilon}$ & \textbf{Calidad} & \textbf{Cota Teórica} & \textbf{Garantía} & \textbf{Estados} & $\mathbf{T}$ (ms) \\
        \midrule
Inst\_01 & 10 & 153 & 0.50 & 566 & 566 & 100.0\% & 283.0 & SI & 1892 & 0.347 \\
Inst\_01 & 10 & 153 & 0.25 & 566 & 566 & 100.0\% & 424.5 & SI & 3839 & 0.646 \\
Inst\_01 & 10 & 153 & 0.10 & 566 & 566 & 100.0\% & 509.4 & SI & 9658 & 1.665 \\
Inst\_01 & 10 & 153 & 0.05 & 566 & 566 & 100.0\% & 537.7 & SI & 19360 & 3.265 \\
Inst\_07 & 16 & 171 & 0.50 & 451 & 450 & 99.8\% & 225.5 & SI & 7106 & 1.275 \\
Inst\_07 & 16 & 171 & 0.25 & 451 & 451 & 100.0\% & 338.2 & SI & 14297 & 2.600 \\
Inst\_07 & 16 & 171 & 0.10 & 451 & 451 & 100.0\% & 405.9 & SI & 35887 & 6.884 \\
Inst\_07 & 16 & 171 & 0.05 & 451 & 451 & 100.0\% & 428.4 & SI & 71944 & 13.221 \\
Inst\_10 & 10 & 147 & 0.50 & 226 & 223 & 98.7\% & 113.0 & SI & 1122 & 0.244 \\
Inst\_10 & 10 & 147 & 0.25 & 226 & 226 & 100.0\% & 169.5 & SI & 2288 & 0.393 \\
Inst\_10 & 10 & 147 & 0.10 & 226 & 226 & 100.0\% & 203.4 & SI & 5742 & 0.963 \\
Inst\_10 & 10 & 147 & 0.05 & 226 & 226 & 100.0\% & 214.7 & SI & 11528 & 1.887 \\
Inst\_12 & 16 & 285 & 0.50 & 639 & 638 & 99.8\% & 319.5 & SI & 5627 & 1.036 \\
Inst\_12 & 16 & 285 & 0.25 & 639 & 638 & 99.8\% & 479.2 & SI & 11373 & 2.042 \\
Inst\_12 & 16 & 285 & 0.10 & 639 & 639 & 100.0\% & 575.1 & SI & 28662 & 9.000 \\
Inst\_12 & 16 & 285 & 0.05 & 639 & 639 & 100.0\% & 607.0 & SI & 57426 & 10.854 \\
Inst\_16 & 11 & 257 & 0.50 & 544 & 544 & 100.0\% & 272.0 & SI & 1704 & 0.404 \\
Inst\_16 & 11 & 257 & 0.25 & 544 & 544 & 100.0\% & 408.0 & SI & 3492 & 0.680 \\
Inst\_16 & 11 & 257 & 0.10 & 544 & 543 & 99.8\% & 489.6 & SI & 8796 & 1.563 \\
Inst\_16 & 11 & 257 & 0.05 & 544 & 544 & 100.0\% & 516.8 & SI & 17628 & 3.734 \\
Inst\_19 & 14 & 189 & 0.50 & 469 & 465 & 99.2\% & 234.5 & SI & 4710 & 1.004 \\
Inst\_19 & 14 & 189 & 0.25 & 469 & 469 & 100.0\% & 351.8 & SI & 9525 & 1.927 \\
Inst\_19 & 14 & 189 & 0.10 & 469 & 469 & 100.0\% & 422.1 & SI & 23880 & 4.432 \\
Inst\_19 & 14 & 189 & 0.05 & 469 & 469 & 100.0\% & 445.6 & SI & 47820 & 9.216 \\
        \bottomrule
    \end{tabular*}
\end{table}

\noindent\textbf{Conclusión de la tabla:} En el $100\%$ de las 80 ejecuciones se satisfizo holgadamente la garantía teórica $\text{ALG}_\varepsilon \ge (1 - \varepsilon)\,\text{OPT}$.

\subsection{Respuesta a la pregunta de análisis: calidad vs costo y mesetas}
\textbf{Pregunta de la guía:} \textit{Explicar por qué reducir $\varepsilon$ normalmente mejora la calidad y aumenta el costo. ¿Pueden existir situaciones donde reducir $\varepsilon$ no mejore la calidad?}

\subsubsection*{1. ¿Por qué mejora la calidad?}
Al redondear $v'_i = \lfloor v_i / K \rfloor$, se pierde como máximo $K$ de valor por objeto. Para el conjunto óptimo de a lo sumo $n$ elementos, la pérdida máxima es:
\begin{equation}
    \text{Pérdida} \le n K = n \cdot \frac{\varepsilon V_{\max}}{n} = \varepsilon V_{\max} \le \varepsilon\,\text{OPT}
\end{equation}
Al reducir $\varepsilon$ (por ejemplo de $0.50$ a $0.05$), el margen máximo de error cae del $50\%$ a tan solo $5\%$, obligando a la solución aproximada a converger al óptimo.

\subsubsection*{2. ¿Por qué aumenta el costo computacional?}
Los valores escalados son $v'_i \le n/\varepsilon$, por lo que la suma de valores escalados es $\sum v'_i \le n^2/\varepsilon$. Las dimensiones de la tabla dinámica son $(n + 1) \times (\sum v'_i + 1)$, de modo que:
\begin{equation}
    \text{Estados procesados y Tiempo} \in O\left( \frac{n^3}{\varepsilon} \right)
\end{equation}
El costo crece de manera \textbf{estrictamente proporcional a $1/\varepsilon$} (crecimiento lineal en $1/\varepsilon$).

\subsubsection*{3. ¿Cuándo reducir $\varepsilon$ NO mejora la calidad?}
\begin{itemize}
    \item \textbf{Mesetas de óptimo prematuro:} En muchas instancias, incluso con $\varepsilon = 0.50$, el orden relativo de los objetos ya reproduce exactamente el óptimo ($\text{ALG}_{0.50} = \text{OPT}$). Reducir $\varepsilon$ a $0.05$ no puede mejorar el valor real ya óptimo, pero sí multiplica el tiempo por diez.
    \item \textbf{Efecto discreto del piso $\lfloor \cdot \rfloor$:} Pequeños cambios en $K$ pueden alterar empates numéricos en la tabla dinámica, produciendo mesetas locales de calidad.
\end{itemize}

\subsection{Análisis gráfico de Knapsack FPTAS}

\begin{figure}[H]
    \centering
    \begin{subfigure}[b]{0.49\textwidth}
        \centering
        \includegraphics[width=\textwidth]{knap_calidad_vs_epsilon.png}
        \caption{Calidad empírica observada frente a $\varepsilon$}
        \label{fig:knap_calidad}
    \end{subfigure}
    \hfill
    \begin{subfigure}[b]{0.49\textwidth}
        \centering
        \includegraphics[width=\textwidth]{knap_costo_vs_inv_epsilon.png}
        \caption{Costo computacional frente a $1/\varepsilon$}
        \label{fig:knap_costo}
    \end{subfigure}
    \caption{Evaluación experimental del esquema FPTAS para Knapsack 0/1}
    \label{fig:knap_completa}
\end{figure}

\subsubsection*{Puntos clave para la sustentación}
\begin{itemize}
    \item En la Figura \ref{fig:knap_completa} (subfigura \ref{fig:knap_calidad}), la calidad media supera ampliamente la cota teórica $(1 - \varepsilon)$ (línea segmentada roja), manteniéndose por encima del $98\%$ en todo momento.
    \item En la subfigura \ref{fig:knap_costo}, el tiempo de ejecución sigue una pendiente lineal estricta respecto a $1/\varepsilon$, confirmando la propiedad fundamental de un algoritmo FPTAS.
\end{itemize}
